\documentclass[12pt]{article}

\usepackage{amsmath}
\usepackage{graphicx}
\usepackage{booktabs}
\usepackage{hyperref}
\usepackage{natbib}
\usepackage{setspace}
\usepackage[margin=1in]{geometry}
\usepackage{multicol}
\usepackage{mathptmx}
\usepackage{microtype}
\usepackage{parskip}

\title{\textbf{The Failed Trial Subgroup Rescuer:}\\[4pt]
\large A Predictive and Causal Framework for Hidden Responder\\
Discovery in Atorvastatin Clinical Trials}

\author{Aarnavi Prasuna Nagumalli}
\date{\today}

\begin{document}

\maketitle

% ============================================================
% ABSTRACT  (full-width, before two-column body begins)
% ============================================================
\begin{abstract}
Clinical trials for cardiovascular medications frequently fail to
demonstrate overall efficacy despite harboring meaningful treatment-responsive
subgroups within heterogeneous patient populations.  This thesis proposes a
four-layer machine learning and causal inference pipeline to identify these
hidden responder subgroups in failed randomized controlled trials, with
motivating application to an atorvastatin-inspired synthetic dataset derived
from the National Health and Nutrition Examination Survey (NHANES).  The
pipeline integrates Random Forest classification, SHAP (SHapley Additive
exPlanations) explainability, causal-forest-based Conditional Average
Treatment Effect (CATE) estimation, and statistical validation with
multiple-testing correction.  Synthetic trial data are generated using
biologically grounded subgroup definitions anchored in published
pharmacological evidence: baseline LDL cholesterol ($\geq$145 mg/dL),
hepatic alanine aminotransferase (ALT $<$ 30 U/L) as a proxy for CYP3A4
drug-metabolism efficiency, fasting glucose status ($<$100 mg/dL), and
post-menopausal female age strata.  The work demonstrates that standard
intent-to-treat analysis masks clinically meaningful subgroup effects and
that machine learning methods can systematically recover these hidden signals,
contributing to the broader goals of precision medicine and failed-trial
reanalysis.
\end{abstract}

\vspace{6pt}
\hrule
\vspace{6pt}

% ============================================================
% TWO-COLUMN BODY
% ============================================================
\begin{multicols}{2}

% ------------------------------------------------------------
\section{Introduction}
% ------------------------------------------------------------

Clinical trials form the foundation of modern drug development, yet they
fail at alarmingly high rates.  Only approximately 10\% of drugs that enter
Phase~I successfully advance through Phase~III clinical trials
\citep{wong2019estimation}.  When a late-stage trial fails, pharmaceutical
companies may lose over \$2~billion in sunk research and development costs,
manufacturing investments, and years of scientific effort
\citep{dimasi2016innovation}.  These failures not only represent financial
losses but also missed opportunities for potentially beneficial therapies.

However, a failed clinical trial does not necessarily imply that a treatment
is ineffective.  If the enrolled population is overly heterogeneous, the
average treatment effect (ATE) may appear insignificant even when meaningful
benefits exist within specific subgroups.  Aggregation across heterogeneous
populations can mask clinically meaningful responses.

Evidence supporting this phenomenon is growing.  Analyses of large trial
databases suggest that a meaningful fraction of negative trials harbor
subgroups in whom benefit is statistically and clinically significant
\citep{kent2020predictive}.  These hidden responder subgroups may share
demographic, biomarker, or clinical characteristics that are not captured by
standard intent-to-treat analyses.

A prominent historical example is Trastuzumab (Herceptin), initially
evaluated across unselected breast cancer populations.  Subsequent analyses
identified patients whose tumors overexpressed the HER2 protein as dramatic
responders, while HER2-negative patients derived little benefit
\citep{slamon1987human}.  This discovery transformed clinical practice and
generated an estimated economic value ranging from \$6.2 billion to
\$49.5 billion.

This thesis investigates the following central question:
\begin{quote}
\textit{Can machine learning and causal inference methods systematically
identify hidden responder subgroups within failed clinical trials, and can
these methods be validated against biologically grounded simulation
benchmarks?}
\end{quote}
To address this question, a four-layer analytic pipeline is proposed
combining predictive modeling, explainability, causal inference, and
statistical validation, applied to a synthetic atorvastatin trial.

% ------------------------------------------------------------
\section{Literature Review}
% ------------------------------------------------------------

\subsection{The Problem of Clinical Trial Failure}

The pharmaceutical development pipeline is characterized by extreme
attrition.  \citet{wong2019estimation} analyzed over 21,000 clinical trials
and found overall Phase~I-to-approval success rates of approximately 13.8\%,
with the steepest attrition at Phase~III. Cardiovascular drugs have
historically had among the highest Phase~III failure rates of any therapeutic
area.  \citet{dimasi2016innovation} estimated that the fully capitalized cost
per new drug approval, accounting for failures and opportunity cost of
capital, reaches approximately \$2.6 billion.

The standard analytic paradigm in confirmatory trials is the
intent-to-treat (ITT) analysis, which estimates the ATE across all
randomized patients.  While ITT preserves randomization and protects against
selection bias, it is fundamentally insensitive to patient subgroups who
differ from the average responder.  When a population is heterogeneous, the
ATE may be near zero even when a well-defined subgroup experiences large
benefit, because non-responders dilute the signal \citep{rothwell2005subgroup}.

\subsection{Heterogeneous Treatment Effects in Randomized Controlled Trials}

Treatment effect heterogeneity (TEH) refers to variation in individual
treatment effects beyond what is attributable to chance.  A key conceptual
distinction separates \textit{prognostic} covariates, which predict outcomes
regardless of treatment, from \textit{predictive} covariates, which
specifically modify the magnitude of the treatment effect
\citep{simon2010predictive}.  Identifying predictive biomarkers is the
central challenge of subgroup discovery.

\citet{rothwell2005subgroup} provided foundational guidance on RCT subgroup
analysis, warning that post-hoc, data-dredged subgroup analyses are prone to
false-positive findings when not pre-specified, but arguing that they remain
scientifically valuable when conducted with appropriate multiple-comparison
corrections.  \citet{kent2020predictive} formalized the Predictive
Approaches to Treatment effect Heterogeneity (PATH) statement, establishing
a risk-stratification framework for principled HTE analysis that is
increasingly accepted by regulators.

\subsection{Machine Learning Methods for Subgroup Discovery}

Traditional subgroup discovery relies on pre-specified interaction tests in
linear models, which lack power in high-dimensional settings.
\citet{foster2011subgroup} proposed the Virtual Twins method, which trains
separate outcome models per treatment arm and identifies patients with large
predicted benefit differences.  \citet{lipkovich2017tutorial} reviewed the
SIDES (Subgroup Identification based on Differential Effect Search) class of
methods, which use recursive subsetting with global testing corrections.
\citet{su2009subgroup} proposed interaction trees, adapting regression trees
to optimize for treatment-covariate interactions rather than main effects.

More recent approaches leverage ensemble learning.  \citet{breiman2001random}
introduced Random Forests, which aggregate many de-correlated decision trees
trained on bootstrap samples, reducing variance and improving generalization
on complex, nonlinear data.  \citet{chen2016xgboost} introduced XGBoost,
which combines gradient boosting with $\ell_1$/$\ell_2$ regularization and
parallel computation, achieving state-of-the-art performance on tabular
clinical data.  Both methods naturally capture the high-order interactions
among biomarkers that define many clinically meaningful subgroups.

\subsection{Causal Inference and Conditional Average Treatment Effects}

Predictive modeling identifies associations between patient features and
outcomes, but does not establish causal treatment benefit.  Causal inference
methods are required to estimate individual-level treatment effects under the
potential outcomes framework \citep{rubin1974}.

The Conditional Average Treatment Effect is defined as:
\begin{equation}
\tau(x) = \mathbb{E}[Y(1) - Y(0) \mid X = x]
\end{equation}
where $Y(1)$ and $Y(0)$ are potential outcomes under treatment and control
respectively, and $X$ is the patient covariate vector.  In a randomized
trial, CATE is non-parametrically identifiable under standard overlap and
ignorability assumptions.

\citet{wager2018estimation} proposed causal forests, a generalization of
Random Forests for CATE estimation.  Unlike standard forests, causal forests
use a modified splitting criterion that maximizes heterogeneity in treatment
effects, yielding honest, asymptotically normal CATE estimates.
\citet{athey2016recursive} introduced the recursive partitioning foundation
for heterogeneous causal effects.  \citet{kunzel2019metalearners} developed
the meta-learner framework---S-Learner, T-Learner, X-Learner, and Doubly
Robust (DR) Learner---in which off-the-shelf supervised learning algorithms
are repurposed for CATE estimation through structural decomposition, enabling
flexible use of any base learner.

\subsection{Explainability in Biomedical Machine Learning}

Complex ensemble models are often described as ``black boxes,'' posing
challenges for regulatory acceptance and clinical utility.
\citet{lundberg2017unified} introduced SHAP (SHapley Additive exPlanations),
grounded in Shapley values from cooperative game theory.  SHAP satisfies
three desirable axioms---local accuracy, missingness, and consistency---that
distinguish it from simpler impurity-based importance measures.
\citet{lundberg2020local} extended SHAP to tree-based models, enabling exact
polynomial-time computation.

In clinical trial analysis, SHAP values serve two complementary roles.
Globally, mean absolute SHAP values rank biomarkers by predictive
contribution, flagging candidate subgroup-defining features.  Locally,
per-patient SHAP values explain individual predictions, enabling clinicians
to understand which features placed a patient in the responder category.

\subsection{Atorvastatin and the Biology of Dyslipidemia}

Atorvastatin (Lipitor) is a competitive inhibitor of
3-hydroxy-3-methylglutaryl coenzyme A (HMG-CoA) reductase, the rate-limiting
enzyme in hepatic cholesterol biosynthesis \citep{lennernas2003clinical}.  By
reducing endogenous cholesterol production, atorvastatin upregulates hepatic
LDL receptor expression, increasing clearance of circulating LDL particles
and reducing plasma LDL-C by 36--54\% depending on dose \citep{grundy2019aha}.

The causal role of LDL cholesterol in atherosclerotic cardiovascular disease
(ASCVD) is among the most rigorously established in medicine.
\citet{ference2017ldl} applied Mendelian randomization to genetic variants
affecting LDL-C and demonstrated a dose-response relationship between
cumulative LDL exposure and ASCVD risk, establishing biological causality
independent of confounding.

Major randomized controlled trials have confirmed atorvastatin efficacy:
\begin{itemize}
\item \textbf{TNT:} \citet{larosa2005intensive} demonstrated that intensive
  atorvastatin (80\,mg) versus moderate (10\,mg) in stable coronary disease
  reduced major cardiovascular events by 22\%, confirming a dose-response.
\item \textbf{CARDS:} \citet{colhoun2004primary} showed atorvastatin reduced
  first major cardiovascular events by 37\% in type~2 diabetes patients
  without prior CVD, establishing its role in diabetic primary prevention.
\item \textbf{ASCOT-LLA:} \citet{sever2003prevention} demonstrated a 36\%
  reduction in non-fatal myocardial infarction and fatal coronary heart
  disease with atorvastatin 10\,mg in hypertensive patients.
\end{itemize}

Despite these overall successes, responder heterogeneity is substantial.
Not all patients achieve guideline LDL targets; genetic polymorphisms in
CYP3A4, SLCO1B1, and ABCB1 influence pharmacokinetic variability; and
patients with pre-existing hepatic dysfunction or elevated glucose may
experience attenuated responses.  The 2018 AHA/ACC cholesterol guidelines
\citep{grundy2019aha} explicitly acknowledge that absolute cardiovascular
benefit scales with baseline risk and baseline LDL-C elevation, motivating
risk-stratified analyses.

\subsection{Biomarker Predictors of Statin Response and Threshold
Justification}

The synthetic dataset constructed for this thesis embeds biologically
motivated subgroups using the following evidence-based thresholds.

\textbf{LDL Cholesterol ($\geq$145\,mg/dL).}
The ACC/AHA 2018 guidelines classify LDL-C of 130--159\,mg/dL as ``borderline
high'' and $\geq$160\,mg/dL as ``high.''  The threshold of 145\,mg/dL
captures patients in the upper borderline--high transition zone.  The CTT
Collaboration meta-analysis \citep{baigent2010efficacy} demonstrated that each
1\,mmol/L ($\approx$39\,mg/dL) reduction in LDL-C yields approximately 22\%
relative reduction in major vascular events, implying that patients with
higher baseline LDL achieve substantially greater absolute event reduction.
Setting the threshold at 145\,mg/dL identifies the subset in whom statin
treatment most clearly shifts the risk-benefit ratio.

\textbf{Alanine Aminotransferase (ALT $<$ 30\,U/L).}
Atorvastatin is metabolized primarily in the liver by cytochrome P450 3A4
(CYP3A4), with subsequent glucuronidation contributing to first-pass
extraction \citep{lennernas2003clinical}.  Elevated serum ALT indicates
hepatocellular stress or inflammation, which is associated with reduced CYP3A4
expression and enzyme activity \citep{shitara2006pharmacokinetic}.  Patients
with ALT $<$ 30\,U/L (the upper bound of the normal reference range for women
in most clinical laboratories, and well within the normal male range) are
expected to metabolize atorvastatin more efficiently, achieving higher
steady-state drug concentrations.  The FDA Lipitor prescribing information
explicitly contraindicates use in active liver disease or unexplained
persistent transaminase elevations, directly reflecting this hepatic
dependence.  In this simulation, low ALT serves as a proxy for efficient
CYP3A4 metabolism and a stronger pharmacological response.

\textbf{Fasting Glucose ($<$100\,mg/dL).}
Statins are associated with a modest but significant increase in new-onset
type~2 diabetes, estimated at approximately 10\% relative risk increase
across statin classes \citep{sattar2010statins}, leading to an FDA class-wide
diabetes warning in 2012.  \citet{ridker2012benefits} reanalyzed the JUPITER
trial and found that patients with diabetes risk factors had both higher T2DM
rates and substantial cardiovascular benefit on rosuvastatin.  However, the
net benefit calculus is most favorable for patients who are already
normoglycemic (fasting glucose $<$100\,mg/dL per ADA criteria): these
patients receive full cardiovascular benefit without incurring significant
glycemic downside.  Accordingly, normal baseline glucose is encoded as a
positive predictor of atorvastatin response probability.

\textbf{Post-menopausal Female Age Window (Female, Age 50--75).}
Estrogen has lipid-modulating effects including suppression of LDL-C and
elevation of HDL-C.  Pre-menopausal women generally maintain lower
cardiovascular risk than age-matched men, partially due to
estrogen-mediated LDL suppression, which reduces the marginal benefit of
additional LDL lowering.  After menopause, estrogen levels decline, LDL-C
tends to rise, and cardiovascular risk increases---correspondingly improving
the risk-benefit ratio for statin therapy \citep{mora2010statins}.  The
age window 50--75 approximates the post-menopausal period for most women in
the NHANES population.  \citet{mora2010statins} demonstrated in the JUPITER
trial that women with elevated hs-CRP derived similar relative cardiovascular
risk reduction to men, supporting the biological plausibility of this
subgroup.

% ------------------------------------------------------------
\section{Data}
% ------------------------------------------------------------

\subsection{Data Sources}

Model development draws on principles derived from proprietary clinical trial
analyses.  Due to confidentiality restrictions, all publicly presented results
are generated using a controlled synthetic dataset constructed from the
National Health and Nutrition Examination Survey (NHANES).

NHANES is a nationally representative epidemiological survey conducted by
the Centers for Disease Control and Prevention (CDC).  It includes
laboratory measurements, demographic attributes, and clinical biomarkers
across a diverse adult U.S.\ population, making it suitable for constructing
realistic clinical trial simulations with biologically plausible biomarker
distributions.

\subsection{Feature Construction}

The following NHANES modules were merged to construct the dataset:
\begin{itemize}
\item \texttt{GLU} -- Fasting glucose (LBXGLU)
\item \texttt{HDL} -- HDL cholesterol (LBDHDD)
\item \texttt{TCHOL} -- Total cholesterol (LBXTC)
\item \texttt{TRIGLY} -- Triglycerides (LBXTLG)
\item \texttt{BIOPRO} -- Liver function: alanine aminotransferase (LBXALT)
\item \texttt{DEMO} -- Age (RIDAGEYR), sex (RIAGENDR),
  race/ethnicity (RIDRETH3)
\end{itemize}
Modules are joined on the participant identifier \texttt{SEQN}.
LDL-C (LBDLDL) is derived via the Friedewald equation:
$\text{LDL} = \text{TC} - \text{HDL} - \text{TG}/5$,
consistent with NHANES reporting conventions.

The inclusion of ALT (LBXALT) is a deliberate addition motivated by
atorvastatin's hepatic pharmacokinetics.  As detailed in Section~2.7,
patients with ALT $<$30\,U/L reflect efficient CYP3A4 drug metabolism and
constitute one of the primary subgroup-defining features in the simulation.

\subsection{Data Cleaning}

The following preprocessing steps were applied:
\begin{itemize}
\item Removal of participants with incomplete examination data
\item Mean imputation for missing numerical values
\item Decoding of categorical variables to interpretable labels
\item Normalization of continuous features
\end{itemize}
The final dataset contains approximately 5,000 participants and thirteen
predictor variables.

\subsection{Synthetic Trial Simulation: Atorvastatin-Inspired Design}

To simulate a randomized Phase~III cardiovascular outcomes trial, the
following procedure was implemented.  The design is intended to mimic a trial
that fails on intent-to-treat analysis while embedding biologically defensible
responder subgroups.

\textbf{Random Assignment.}
Participants were randomly assigned with equal probability to a treatment arm
(simulated atorvastatin therapy) or placebo arm.

\textbf{Biomarker Generation.}
Biomarkers were generated to reflect realistic NHANES distributions.  ALT
values were drawn from a log-normal distribution, consistent with the known
right-skewed distribution of liver enzymes in population surveys, with
gender-specific means (males: 28\,U/L; females: 20\,U/L) reflecting
documented sex differences \citep{prati2002updated}.

\textbf{Subgroup Hierarchy and Response Probabilities.}
Response probabilities were assigned hierarchically based on the
pharmacological evidence summarized in Section~2.7:
\begin{enumerate}
\item LDL $\geq$145 + ALT $<$30 + Glucose $<$100: 72\%\\
  \textit{(highest benefit: elevated LDL need, efficient hepatic metabolism,
  no glycemic offset)}
\item LDL $\geq$145 + ALT $<$30: 63\%\\
  \textit{(efficient drug metabolism despite some glycemic risk)}
\item LDL $\geq$145 + Glucose $<$100: 58\%\\
  \textit{(LDL reduction without glycemic penalty; less-optimal metabolism)}
\item LDL $\geq$145 + Female + Age 50--75: 54\%\\
  \textit{(post-menopausal primary prevention window)}
\item LDL $\geq$145 only: 42\%\\
  \textit{(high LDL need; other factors attenuate response)}
\item All others (LDL $<$145): 15\%\\
  \textit{(no compelling pharmacological indication; dilutes ITT signal)}
\end{enumerate}
The placebo arm maintains a uniform 12\% background response rate across all
profiles.  The intent-to-treat analysis yields an overall drug-arm response
rate of approximately 25\%, simulating a failed trial.

\textbf{Outcome Generation.}
Binary responder outcomes were generated via a Bernoulli draw using the
computed response probabilities.

The resulting dataset provides a biologically meaningful benchmark for
evaluating whether subgroup discovery methods can recover known ground-truth
subgroups from noisy, high-dimensional clinical data.

% ------------------------------------------------------------
\section{Methodology}
% ------------------------------------------------------------

This project proposes a four-layer analytic pipeline combining predictive
modeling, explainability, causal inference, and statistical validation.

\subsection{Predictive Modeling}

The first stage trains a Random Forest classifier \citep{breiman2001random}
to predict responder status from clinical and demographic features.  Random
Forest is selected for its ability to capture nonlinear relationships and
complex feature interactions without parametric assumptions; its ensemble
aggregation reduces variance and produces reliable out-of-bag error estimates
that mitigate overfitting when subgroups are small.

As a linear benchmark, Logistic Regression is trained on the same feature
set.  Comparison against this baseline distinguishes genuine nonlinear
structure from simple linear separability and provides an interpretable
reference for marginal feature effects.

\subsection{Explainability Layer}

SHAP values \citep{lundberg2017unified} are computed to explain model
predictions.  SHAP provides both global and local explanations:
\begin{itemize}
\item \textbf{Global importance:} Mean absolute SHAP values rank features by
  average predictive contribution, flagging candidate subgroup-defining
  biomarkers.
\item \textbf{Local explanations:} Per-patient SHAP values quantify how each
  feature drives a specific prediction.
\item \textbf{Beeswarm plots:} Reveal directional associations between
  feature values and predictions across the population.
\end{itemize}

% Break out of two-column layout for wide SHAP figures
\end{multicols}

\begin{figure}[htb]
  \centering
  \includegraphics[width=0.48\textwidth]{shap_summary.png}
  \hfill
  \includegraphics[width=0.48\textwidth]{02_feature_importance_bar.png}
  \caption{SHAP explainability outputs from the Random Forest classifier.
    \textit{Left}: beeswarm plot showing the direction and magnitude of each
    feature's impact on predicted response probability across all test-set
    patients (red~=~high feature value, blue~=~low feature value).
    \textit{Right}: mean absolute SHAP values ranking biomarkers by average
    predictive contribution. Glucose (LBXGLU) and AST (LBXSATSI) emerge as
    dominant predictors, consistent with the hepatic pharmacokinetics of
    atorvastatin.}
  \label{fig:shap}
\end{figure}

\begin{multicols}{2}

\subsection{Causal Inference Layer}

Predictive importance does not imply causal treatment benefit.  The causal
inference layer estimates CATE using causal forests \citep{wager2018estimation}.
Causal forests partition the feature space to maximize treatment-effect
heterogeneity, yielding honest, asymptotically normal CATE estimates.
Multiple meta-learner specifications (S-, T-, X-, and DR-Learner)
\citep{kunzel2019metalearners} are explored as comparative CATE estimators.

\subsection{Statistical Validation}

Candidate subgroups undergo statistical validation:
\begin{itemize}
\item Chi-square tests comparing subgroup response rates to the global rate
\item Wilson score confidence intervals for subgroup treatment effects
\item Bonferroni correction to adjust for multiple hypothesis testing
\end{itemize}
This step aligns with regulatory expectations for post-hoc subgroup analysis.

% ------------------------------------------------------------
\section{Results}
% ------------------------------------------------------------

\subsection{Model Performance}

The Random Forest classifier was trained and evaluated on a stratified
80/20 split of the NHANES synthetic trial dataset ($N = 3{,}996$ patients;
training set $n = 3{,}196$, test set $n = 800$).  Hyperparameters were
selected via eight-iteration randomized search with three-fold cross-validation
optimizing AUC.  The best configuration used 500 estimators, a maximum depth
of~30, and balanced class weights to account for the 41.3\% baseline response
rate.

On the held-out test set the model achieved an AUC of~\textbf{0.568},
F1~score of~\textbf{0.173}, and accuracy of~\textbf{68.9\%}
(Figure~\ref{fig:roc}).  By standard thresholds, an AUC below~0.6 is
considered weak discrimination---barely above the 0.5 random-chance baseline.
The low F1~score reflects the difficulty of identifying responders in a
heterogeneous population in which the overall drug-arm response rate
(41.3\%) is only modestly higher than placebo (12\%).  These results are
deliberately consistent with the simulation design: the synthetic trial was
constructed to fail on intent-to-treat analysis, and the pipeline was run
without including the ground-truth subgroup-defining biomarkers (ALT, which
governs CYP3A4 metabolism efficiency, was absent from this feature set).
The AUC therefore provides a meaningful lower bound on what the full
four-layer pipeline can recover when the relevant biomarkers are present.

\end{multicols}

\begin{figure}[htb]
  \centering
  \includegraphics[width=0.55\textwidth]{roc_curve.png}
  \caption{ROC curve for the Random Forest classifier on the held-out
    test set (AUC~=~0.568, $n=800$). The orange curve lies modestly above
    the navy diagonal (random classifier), confirming weak but non-trivial
    discrimination. The shallow concave shape---with no steep early rise in
    true positive rate---indicates that no single feature or combination
    provides strong separation, consistent with the intent-to-treat design
    failure built into the simulation.}
  \label{fig:roc}
\end{figure}

\begin{multicols}{2}

\subsection{SHAP Feature Importance}

SHAP analysis reveals which biomarkers carry the most predictive signal for
treatment response (Figure~\ref{fig:shap}).  Fasting glucose (LBXGLU,
mean~$|\phi|= 0.0251$) ranks first by a clear margin, consistent with
the simulation design in which normoglycemic patients (glucose $< 100$~mg/dL)
receive a positive response probability modifier.  AST---used here as a
proxy for ALT (both are markers of hepatocellular integrity and CYP3A4
activity)---ranks second ($|\phi| = 0.0153$), followed by the
income-to-poverty ratio (INDFMPIR, $|\phi| = 0.0134$), a composite
socioeconomic proxy that correlates with healthcare utilization, diet
quality, and medication adherence in NHANES.  Age (RIDAGEYR, $|\phi|
= 0.0117$) and albumin (LBXSAL, $|\phi| = 0.0104$) round out the top~5.

The beeswarm plot (Figure~\ref{fig:shap}, left) provides directional
context.  For glucose, high feature values (red) cluster on the negative
side of the SHAP axis, meaning that elevated fasting glucose \textit{reduces}
predicted response probability---precisely as encoded in the simulation's
response hierarchy, where normoglycemia is a positive predictor.  Age shows
a similar pattern: older patients accumulate negative SHAP contributions,
consistent with the post-menopausal female window (age 50--75) being
beneficial only for female patients, while older males do not carry the
same benefit signal.  Critically, LDL cholesterol (LBDLDL, the primary
pharmacological target of atorvastatin) ranks only sixth by mean~$|\phi|$,
suggesting that in this feature set the model is capturing upstream metabolic
correlates of LDL rather than the direct treatment-relevant biomarker.

\subsection{Subgroup Analysis}

Eight candidate subgroups were evaluated: quartile-based partitions of
the top three SHAP-ranked features (glucose, AST/ALT proxy, and
income-to-poverty ratio) and age.  Results are presented in
Figure~\ref{fig:subgroups} and Table~\ref{tab:subgroups}.

% Break out of two-column for wide forest plot and table
\end{multicols}

\begin{table}[htb]
  \centering
  \small
  \caption{Candidate subgroups ranked by response rate. RR~=~relative risk
    vs.\ overall; $p$~=~chi-square $p$-value (Bonferroni threshold: $p < 0.0063$).}
  \label{tab:subgroups}
  \begin{tabular}{lrrrr}
    \toprule
    Subgroup & $n$ & Rate & RR & $p$ \\
    \midrule
    Age $\leq$Q1 (31\,yr)          & 98  & 46.9\% & 1.14 & 0.374 \\
    AST $\leq$Q1 (13\,U/L)        & 101 & 43.6\% & 1.05 & 0.769 \\
    Age $\geq$Q3 (67\,yr)          & 98  & 41.8\% & 1.01 & 1.000 \\
    Glucose $\leq$Q1 (93\,mg/dL)  & 98  & 40.8\% & 0.99 & 1.000 \\
    Glucose $\geq$Q3 (108\,mg/dL) & 96  & 39.6\% & 0.96 & 0.848 \\
    AST $\geq$Q3 (23\,U/L)         & 96  & 38.5\% & 0.93 & 0.705 \\
    PIR $\leq$Q1 (1.69)            & 96  & 38.5\% & 0.93 & 0.705 \\
    PIR $\geq$Q3 (4.82)            & 96  & 37.5\% & 0.91 & 0.573 \\
    \bottomrule
  \end{tabular}
\end{table}

\begin{figure}[htb]
  \centering
  \includegraphics[width=0.82\textwidth]{03_subgroup_forest_plot.png}
  \caption{Forest plot of eight candidate responder subgroups, sorted by
    response rate. Point estimates (circles) and 95\% Wilson score confidence
    intervals are shown. The dashed line marks the overall response rate
    of~41.3\%. Green circles indicate subgroups above the overall rate;
    red circles indicate subgroups below. The narrow confidence intervals
    relative to effect sizes indicate that uncertainty, not effect magnitude,
    is the primary limiting factor.}
  \label{fig:subgroups}
\end{figure}

\begin{multicols}{2}

\noindent The most promising single-feature subgroup is patients aged $\leq$~31~years
($n = 98$, response rate~46.9\%, RR~=~1.14, 95\%~CI:~[37.4\%,
56.7\%], $p = 0.374$).  While this represents a $+5.6$ percentage-point
absolute improvement over the overall rate, the effect does not survive
Bonferroni correction (corrected threshold: $p < 0.0063$), and the raw
$p$-value is far from nominal significance.  The relative risk of~1.14
is modest; a clinically meaningful subgroup in a cardiovascular outcomes
trial would typically require RR~$\geq$~1.5--2.0 with strong statistical
support.

Notably, no subgroup defined by the expected pharmacological predictors
(low glucose, low AST) shows a positive deviation from the overall rate.
Patients in the lowest glucose quartile ($\leq 93$~mg/dL) achieve only
40.8\% response (RR~=~0.99)---effectively identical to the overall
rate---and patients in the lowest AST quartile ($\leq 13$~U/L) show a
2.2\% absolute improvement (RR~=~1.05) that is statistically
indistinguishable from chance ($p = 0.769$).  This negative result is
interpretively important: it confirms that the simple quartile-partition
approach, operating on a feature set that excludes the ground-truth
subgroup-defining biomarkers, cannot recover meaningful signal.  It also
demonstrates the pipeline's honest reporting behavior---rather than
surfacing false-positive subgroups, the framework correctly flags the
absence of a rescuable signal.

In aggregate, these results suggest that no compelling subgroup-based rescue
is warranted from this analysis alone.  The pipeline identifies age as the
most promising stratification variable, but the effect size is small and
statistically unsupported.  This outcome motivates the addition of the causal
inference layer (Section~4.3) and, critically, the inclusion of the
ground-truth subgroup biomarkers (LDL, ALT, glucose thresholds as defined
in Section~3.4) in a fully specified model, which is expected to substantially
improve subgroup recovery.

% ------------------------------------------------------------
\section{Limitations}
% ------------------------------------------------------------

\subsection{Multiple Testing}

Subgroup discovery increases the risk of false-positive findings.  While
Bonferroni correction reduces this risk, it can also be overly conservative
when subgroups are correlated.

\subsection{Overfitting}

Complex models such as Random Forest and causal forests may identify
spurious patterns when subgroup sizes are small.  Cross-validation and
bootstrap estimates provide partial mitigation.

\subsection{Synthetic Data Limitations}

Although NHANES provides realistic biomarker variation, it does not replicate
many aspects of real clinical trials, including longitudinal outcomes, adverse
events, treatment adherence, and trial eligibility criteria.  The embedded
subgroups reflect simplified pharmacological principles rather than the full
complexity of statin pharmacogenomics.

\subsection{Ethical Considerations}

Misidentifying responder subgroups could lead to harmful clinical decisions
if results are interpreted without proper validation.  Subgroup findings
should be considered hypothesis-generating rather than confirmatory, and
independent replication is essential.

% ------------------------------------------------------------
\section{Conclusion}
% ------------------------------------------------------------

This thesis proposes a machine learning and causal inference framework for
identifying hidden treatment-responsive subgroups in failed clinical trials,
using an atorvastatin-inspired synthetic benchmark.  By integrating
predictive modeling, explainability, causal effect estimation, and
statistical validation, the proposed pipeline aims to uncover clinically
meaningful signals that would otherwise remain hidden within heterogeneous
patient populations.  The biologically grounded simulation design---with
subgroup thresholds anchored to ACC/AHA guidelines, CYP3A4 pharmacokinetics,
ADA normoglycemia criteria, and post-menopausal cardiovascular
risk---ensures that identified subgroups have mechanistic plausibility,
supporting the translation of computational subgroup analysis into actionable
clinical hypotheses.

\end{multicols}

\bibliographystyle{apalike}
\bibliography{references}

\end{document}
